\chapter*{Introduzione al Corso e Modalità d'Esame}

\section*{Descrizione dell'Insegnamento}
L'insegnamento di Data Science e Tecnologie per le Basi di Dati è un corso obbligatorio per la Laurea Magistrale in Ingegneria Informatica. L'obiettivo principale del corso è descrivere gli aspetti fondamentali della tecnologia dei sistemi per la gestione delle basi di dati (DBMS) e le metodologie di gestione delle basi di dati destinate all'analisi (comunemente denominate data warehouse). 

Questi sistemi si caratterizzano per la necessità di gestire enormi volumi di dati. Nel corso vengono esplorate sia le tecniche di analisi tradizionali, come l'analisi OLAP (On Line Analytical Processing), sia tecniche di analisi più complesse basate sul data mining.

\section*{Programma e Argomenti Trattati}
Il programma del corso si articola in diverse macro-aree fondamentali per la gestione e l'analisi dei dati:
\begin{itemize}
    \item \textbf{Caratteristiche tecnologiche dei DBMS:} vengono studiati la gestione dell'accesso concorrente ai dati, l'affidabilità del sistema, le strutture del livello fisico e l'ottimizzazione dell'accesso ai dati.
    \item \textbf{Basi di dati NoSQL:} introduzione al modello dei dati e ai linguaggi di interrogazione di sistemi non relazionali, con particolare riferimento a MongoDB ed Elastic.
    \item \textbf{Basi di dati distribuite:} studio dell'architettura dei sistemi distribuiti e delle relative tecniche di gestione della replicazione.
    \item \textbf{Data warehouse (DW):} analisi dell'architettura e delle metodologie di progettazione concettuale, logica e fisica. Vengono inoltre trattate la preparazione dei dati e la formulazione di interrogazioni OLAP tramite il linguaggio SQL.
    \item \textbf{Data science e Machine learning:} panoramica sul processo di data science e approfondimento dei principali algoritmi di apprendimento automatico per la classificazione, il clustering e l'estrazione di regole di associazione.
\end{itemize}

\section*{Struttura e Regole dell'Esame}
L'esame si compone di una prova scritta obbligatoria e della valutazione di eventuali homework assegnati durante il semestre (elaborato progettuale individuale). Gli homework sono facoltativi, ma fortemente consigliati. A discrezione del docente, potrebbe essere richiesta una prova integrativa a conferma del voto ottenuto.

\subsection*{La Prova Scritta}
La prova scritta si svolge in aula tramite PC, utilizzando la piattaforma di ateneo, e ha una durata complessiva di 100 minuti. Durante lo svolgimento non è permesso l'utilizzo di libri, appunti o dispositivi elettronici di alcun tipo.

Il test è composto da domande a risposta aperta e chiusa, oltre a esercizi pratici di progettazione. Le regole di valutazione delle domande prevedono che le risposte errate ai quesiti a risposta chiusa comportino una penalizzazione, mentre le risposte omesse valgano zero.

L*a struttura tipica della prova scritta è la seguente:
\begin{itemize}
    \item \textbf{Domande teoriche (6-8 quesiti):} vertono sui concetti principali del corso, quali aspetti tecnologici dei DBMS, basi di dati distribuite, preparazione dei dati e algoritmi di analisi (massimo 8 punti).
    \item \textbf{Progettazione fisica (1-3 esercizi):} esercizi sulla creazione di strutture fisiche per basi di dati (massimo 3 punti).
    \item \textbf{Progettazione Data Warehouse (1-2 esercizi):} modellazione di un DW (massimo 5 punti).
    \item \textbf{Preparazione dei dati (1 esercizio):} operazioni di data preparation per data warehouse in linguaggio SQL (massimo 5 punti).
    \item \textbf{Interrogazioni OLAP (3 esercizi):} scrittura di istruzioni SQL per l'accesso e l'analisi dei dati di un data warehouse (massimo 11 punti).
\end{itemize}

\subsection*{Homework e Valutazione Finale}
Durante il corso verranno proposti quattro homework individuali. Questi lavori sono mirati a sviluppare il pensiero critico e le capacità di problem solving sugli argomenti trattati a lezione. 

Per quanto riguarda il calcolo del voto finale:
\begin{itemize}
    \item Gli homework possono conferire un massimo di 2 punti aggiuntivi.
    \item Le relazioni degli homework vengono prese in considerazione per il voto finale esclusivamente se la valutazione della prova scritta è pari o superiore a 18/30.
    \item La lode viene assegnata nel caso in cui il punteggio complessivo superi il 31/30.
\end{itemize}